\section{Preliminary Feasibility Probe}
\label{sec:pilot}

% =========================================================================
% This section was formerly a results-bearing pilot. It has been neutralized:
% the earlier pilot ran against an action space (and a set of LLM substrates
% and search budgets) that predate the current verb-set redesign of
% Section~\ref{sec:method-actions}, so its specific numbers no longer describe
% the system this paper specifies. We retain it only as a feasibility probe:
% an account of what we built and exercised end-to-end, and of the
% methodological lessons that survive the redesign and feed the pre-registered
% plan of Section~\ref{sec:experimental-plan}. No accuracy or agreement number
% from the earlier runs is reported as evidence, because those runs are
% superseded.
% =========================================================================

This section reports a preliminary feasibility probe conducted on an early
implementation of the LLM context MDP scaffold of
Section~\ref{sec:method}. Its purpose is narrow, and it makes no empirical
claim about benchmark performance: it records that every component of the
framework can be implemented and wired together end-to-end, and it surfaces a
small number of methodological lessons that shape the pre-registered study of
Section~\ref{sec:experimental-plan}. The probe predates the current
action-space design (Section~\ref{sec:method-actions}) and ran on substrates
and search budgets below those the pre-registered plan mandates. We therefore
treat it as superseded scaffolding rather than as a result: it is the
integration check whose purpose was to falsify the framework's basic mechanics
before committing to the benchmark study, not a measurement that the framework's
claims rest on. The empirical claim of the program is the pre-registered
benchmark study of Section~\ref{sec:experimental-plan}, and that study is the
only place this paper reports headline numbers.

\subsection{What the Probe Exercised}
\label{sec:pilot-setup}

The probe instantiated the full loop on a handful of toy problems spanning a few
families of verb-choice difficulty (single-step arithmetic, single-step word
problems, two-step linear algebra, and quadratic factoring), with a second,
structurally parallel set of problems held out and never used during fitting.
The held-out construction is the one design element worth carrying forward
verbatim: held-out problems share structural shape but never surface content
with the training problems, so a model that has merely memorized verb sequences
on training problems cannot score well on the held-out set unless those
sequences generalize. The pre-registered benchmark study inherits this
train/held-out discipline at scale (Section~\ref{sec:experimental-plan}).

The base language model was a small instruction-tuned model served locally,
acting both as the verb instantiator and as the state encoder; the agent's
trainable component was the small multi-head architecture of
Section~\ref{sec:method-architecture}, well under the parameter budget the plan
allows. Search ran at a low simulation budget, far below the pre-registered
$N_{\mathrm{small}} = 50$ and $N_{\mathrm{large}} = 500$
(Section~\ref{sec:exp-frontier}), and the verifier was the SymPy equality call
on canonicalized normal forms of Section~\ref{sec:method-reward}. We mention the
substrate, budget, and action space only to mark how far the probe sits from the
pre-registered configuration: it used a different (now-superseded) verb set, a
larger and instruction-tuned base model than the base-only substrate the plan
mandates for clean attribution (Section~\ref{sec:exp-models}), and a search
budget at the low end of the range the plan commits to. None of the probe's
quantitative outcomes are reported here as evidence for the framework, precisely
because each of these axes changed in the redesign.

\subsection{Round-Over-Round Improvement: Feasibility, Not a Result}
\label{sec:pilot-v2-v3}

The probe collected MCTS visit-count targets under a uniform PUCT prior,
distilled them into the parametric heads, re-collected targets with the
distilled heads installed as the PUCT prior, and re-distilled, closing one round
of expert iteration offline. The qualitative observation that survives the
redesign is the one the framework predicts: distilling sharper MCTS visit-count
targets into the parametric heads is implementable and the offline distillation
step converges, so the policy-improvement-then-distill loop of
Section~\ref{sec:method-training} can be run end-to-end. We do not report the
earlier round-over-round agreement or loss numbers, nor a figure built on them,
because they were measured against the superseded action space and a single
training seed, below the plan's three-seed floor (Section~\ref{sec:exp-stats}).
The single feasibility takeaway is that the loop closes mechanically; whether it
improves accuracy on standardized benchmarks under the pre-registered protocol
is the open empirical question, deferred to Section~\ref{sec:experimental-plan}.

\subsection{A Naive Second-Round Expansion and Its Lesson}
\label{sec:pilot-v3-v4}

The probe also ran a second round of expert iteration on a naively-expanded
training set, adding problems of the family the first-round heads handled worst.
The episode is worth recording for its methodological content rather than its
numbers: expanding the training set of one family while leaving the others fixed
shifts the per-class composition, and a multi-round expansion can change
held-out behavior on families that were not the target of the intervention. The
lesson that carries forward is a monitoring requirement, not a measured effect:
any multi-round result in the benchmark study must report the per-round
training-set composition and the per-class held-out accuracy, so that a
round-over-round change is attributed to its actual mechanism rather than to the
most visible covariate. We state this as a forward-looking commitment for the
pre-registered study (Section~\ref{sec:pilot-lessons}), not as a finding the
probe's superseded data establishes.

\subsection{Disentangling Composition From Stability}
\label{sec:pilot-v4b}

The probe attempted to separate two candidate explanations for round-over-round
regression on the hardest family: class composition (the relative training share
of that family) versus target stability (how reproducibly the MCTS visit-count
target lands on the same verb across independent collection runs). The
qualitative conclusion that survives is that composition correction alone did
not account for the regression, which redirected attention to the stability of
the distillation target at the low search budget the probe used. This is the
hypothesis the budget--stability probe of Section~\ref{sec:pilot-bsweep} was
built to test, and it is the reason the pre-registered study commits to a
target-stability diagnostic alongside any multi-round comparison
(Section~\ref{sec:pilot-lessons}). We report the redirection of attention, not
the specific agreement or distance figures that prompted it, because those
figures belong to the superseded runs.

\subsection{What Stabilizes the Distillation Target}
\label{sec:pilot-bsweep}

The most transferable lesson from the probe concerns how to \emph{measure} the
stability of the MCTS distillation target, and how that stability scales with
the search budget. For a fixed initial state we ran many independent MCTS
searches (independent random seeds, hence independent stochastic LLM
transitions) at a range of simulation budgets on a low-ambiguity control problem
and a high-ambiguity problem, holding the start state fixed so that any
disagreement among the resulting visit-count distributions is the estimator's
own sampling noise rather than a divergent trajectory. Three qualitative
findings survive the redesign and feed the plan.

First, \emph{run-to-run agreement on the recommended verb rises with the search
budget}: at small budgets two independent collections frequently disagree on the
modal verb, and agreement increases as the budget grows, with the low-ambiguity
family converging faster than the high-ambiguity one. The operational
consequence is that the heavy-search budget must be set by the
slowest-converging family rather than the average, which is why the
pre-registered $N_{\mathrm{large}} = 500$ sits well above the budget the probe
collected with.

Second, \emph{distributional distance is the wrong stability metric and pairwise
recommended-verb agreement is the right one}. Mean $L_1$ between two visit-count
distributions is confounded by sharpness: as the budget grows the distributions
sharpen, and a single sharp minority run on a different verb inflates $L_1$ even
as true run-to-run agreement improves, so $L_1$ can anti-correlate with the
quantity it is meant to track. We therefore pre-register \emph{pairwise
recommended-verb argmax agreement across independent searches from a fixed
state} as the target-stability diagnostic (Section~\ref{sec:pilot-lessons}),
reported stratified by problem family rather than pooled.

Third, \emph{small-sample stability diagnostics can manufacture qualitative
structure}. A coarse probe with few searches per cell produced an apparent
non-monotonicity in the budget--stability curve that dissolved once the number
of independent searches per cell was raised. Conclusions about the shape of the
budget--stability curve need a generous number of independent searches per cell,
well above what a quick probe suggests. We report the procedure and these
qualitative lessons; the per-cell numeric table that an early sweep produced is
superseded and is not carried into this paper.

\subsection{Closing the Loop End-to-End}
\label{sec:pilot-closedloop}

The probe also exercised the full online loop rather than a single offline
distillation pass: search with the current heads as the PUCT prior and the
target-$Q$ head as the bootstrap leaf value, distill the resulting $\pimcts$ and
$\Qmcts$ into the heads, then re-prior the next round of search with the updated
heads (Section~\ref{sec:method-training}). Two pieces of infrastructure that
this exercise forced are load-bearing for the pre-registered study and survive
the redesign.

\paragraph{In-process frozen model and honest per-decision compute accounting.}
An HTTP inference endpoint can neither expose the base model's hidden state,
which the adapter architecture of Section~\ref{sec:method-architecture}
requires, nor make the decoder forwards countable for compute accounting. We
therefore moved the frozen model in process and made a single state prefill
serve both the heads and verb generation (``one forward, not two''), so that the
decoder forwards executed per decision become countable. The qualitative finding
that carries forward, independent of the superseded numbers, is that
\emph{generation, not prefill, dominates a decision's compute}: each decision
spends many verb-instantiator continuations producing action content against a
single short state prefill. This is the measurement on which the FLOP-accounted
compute axis of the pre-registered plan (Section~\ref{sec:exp-compute}) is
built: the accounting must charge every generated token inside the search, not
only the root forward, because a simulation's cost is itself dominated by
generation.

\paragraph{Observability and resumability are research infrastructure.}
The framework's central claim is an across-iteration trend, namely the
approximate-$Q$-learning decrease of the $Q$-loss $L_Q$ and the
expert-iteration sharpening of the policy prior $\ptheta$ and the search visits.
An early long run that persisted only a per-iteration $L_{\mathrm{total}}$
summary could not surface that trend, because $L_{\mathrm{total}}$ is dominated
by the policy cross-entropy $L_{\pi}$ sitting near its target-entropy floor,
masking the small $L_Q$ term and the prior-sharpening dynamic. We rebuilt the
run harness so a long run is both observable and recoverable: it records, per
iteration, the separated $L_Q$ and $L_{\pi}$, the Shannon entropy of the visit
distribution $\pimcts$ and of the learned prior $\ptheta$ (the sharpening
diagnostic), and the per-decision FLOP counters, and it checkpoints often enough
that an interrupted run resumes rather than restarts. The pre-registered study
inherits this monitoring by construction rather than retrofitting it onto a long
run. The across-iteration learning signals themselves are not reported here:
they are an empirical claim, deferred to the pre-registered study, and the
earlier runs that touched them used the superseded action space.

\subsection{Lessons for the Pre-Registered Study}
\label{sec:pilot-lessons}

Four commitments carry from the probe into the pre-registered study of
Section~\ref{sec:experimental-plan}. None of them is a result; each is a
methodological precaution the probe surfaced.

First, \emph{train/held-out discipline}: held-out problems share structural
shape but never surface content with the training problems, so generalization
cannot be faked by memorizing verb sequences.

Second, \emph{composition tracking}: any multi-round result reports the
per-round training-set composition and the per-class held-out accuracy, so a
round-over-round change is attributable to mechanism rather than to the most
visible covariate.

Third, \emph{target-stability monitoring}: the pre-registered target-stability
diagnostic is pairwise recommended-verb argmax agreement across independent
searches from a fixed state, reported stratified by problem family rather than
pooled, with the heavy-search budget set by the slowest-converging family.
Agreement below a pre-committed threshold at the chosen budget is a precondition
for trusting any round-over-round improvement claim.
\textcolor{red}{\textbf{[PRE-REG TODO: commit the numerical
target-stability threshold; suggested: pairwise recommended-verb
agreement $\geq 0.70$ at the collection budget, stratified by problem
family, confirm]}}

Fourth, \emph{honest per-decision compute accounting with observability}: the
compute axis charges every generated token inside the search, and the run
harness persists the separated $L_Q$ and $L_{\pi}$, the visit- and
prior-distribution entropies, and the per-decision FLOP counters per iteration,
checkpointing often enough to resume from a crash. The benchmark study inherits
this monitoring by construction.

\subsection{Caveats}
\label{sec:pilot-caveats}

The probe is not a benchmark study and supplies no accuracy claim. It predates
the action-space redesign of Section~\ref{sec:method-actions}, ran on a small
instruction-tuned base model rather than the base-only substrate of
Section~\ref{sec:exp-models}, used a search budget well below the pre-registered
$N_{\mathrm{small}} = 50$ and $N_{\mathrm{large}} = 500$, and fitted heads on a
handful of toy problems with a single training seed against the plan's
three-seed floor (Section~\ref{sec:exp-stats}). For all of these reasons its
quantitative outcomes are superseded and are deliberately not reported. What the
probe establishes is feasibility: every component of Section~\ref{sec:method}
can be implemented and the loop, when implemented, runs end-to-end. Whether the
framework's compute-versus-accuracy frontier clears the pre-registered
falsification bar is the question the benchmark study of
Section~\ref{sec:experimental-plan} exists to answer.

\subsection{Reproducibility}
\label{sec:pilot-reproducibility}

The probe's scaffolding is released alongside this paper for independent
inspection of the implementation, with the explicit caveat that the released
collection and head-training scripts, the in-process model backend, the shared
single-forward encoder, the prioritized replay buffer and target-network sync,
the outer-loop orchestration, and the observable, resumable run harness exercise
an earlier action space than the one this paper specifies. We release them as
evidence that the framework is implementable end-to-end and as the starting
point for the pre-registered benchmark study, not as artifacts whose numeric
outputs support a claim. The pre-registered study's released artifacts
(Section~\ref{sec:exp-stats}) are the ones against which the framework's
empirical claims will be evaluated.
